\section{Parallel Algorithm}
\label{sec:alg}

This section describes our bulk-synchronous-parallel approach to computing
congruence closure. We first describe the two workhorses of our algorithm: a
concurrent union-find data structure used to maintain congruence closure, and
\textit{semisort}, a parallel primitive offered by
\texttt{parlaylib}~\cite{parlaylib} that we use to group congruent terms. We
then describe how we use these components to develop our final algorithm.

\subsection{Concurrent Union-Find}
\label{subsec:uf}

We employ a simple concurrent union-find presented by Alistarh et al.~\cite{uf}.
In their experiments, they study various different union-find algorithms, and
find that most variants perform comparably across their evaluation. In
particular, more theoretically efficient variants (like those employing
randomized linking~\cite{jayanti16}) do not achieve better practical
performance. We therefore chose a simple variant with confidence that it would
perform well in practice, even under contention.

The algorithm employs best-effort linking by rank. The structure is represented
as an array \textit{UF}, where $\mathit{UF}[i]$ is the index of the parent of
node $i$ if $i$ is not a root, or its rank otherwise. The high-order bit
distinguishes whether the entry is another node or the rank of a root. We now
describe the two operations it supports: \textsc{Find} and \textsc{Union}:

\textbf{Find.} A $\Call{Find}{u}$ returns a pair $(\mathit{root},
\mathit{rank})$ comprising the representative of the equivalence class of $u$
and its rank. It first checks whether $\mathit{UF}[u]$ is a rank; if so, it
returns a pair comprising $u$ and that rank. Otherwise, $\mathit{UF}[u]$ is its
parent, and it invokes \textsc{Find} on the parent. It then attempts to path
compress via a \texttt{CAS}, which may fail if the path has already been
compressed.

\textbf{Union.} A $\Call{Union}{u, v}$ operation unions the equivalence classes
of $u$ and $v$. It first computes the roots and ranks of $u$ and $v$ via
\textsc{Find}. It then checks whether the representatives are equal; if so, $u$
and $v$ are already part of the same equivalence class and it returns
immediately. Otherwise, it compares the ranks of $u$ and $v$. If one rank is
strictly greater then the other, then the algorithm attempts to set the parent
of the lower-rank node to the greater-rank node with a \texttt{CAS}. This may
fail if the parent of the lower-rank node has been updated in the meantime, in
which case the operation retries.

If both ranks are equal, then it selects one node arbitrarily to be the parent
and attempts to update the child's parent pointer via a \texttt{CAS}, again
retrying if it fails. It then attempts to increment the rank of the node chosen
to be the parent, which again may fail if another \textsc{Union} operation
interferes. If it fails, it means that the node chosen to be a parent is no
longer a root, and thus has no rank. Thus incrementing the rank can fail, and
linking by rank is best-effort.

\subsection{Parallel Semisort}
\label{subsec:semisort}

At each round, we want to group congruent terms based on newly discovered
equalities. We use the \emph{semisort} of Gu, Shun, Sun, and
Blelloch~\cite{semisort}. Given an array of (hash, term) pairs, we want to sort
the array such that all pairs with the same hash are adjacent. Concretely, we
use semisort to group congruence classes at each round.

\subsection{Rounds-Based Congruence Closure}
\label{subsec:rounds}

At each round, we maintain a set of terms whose congruence classes may need to
be updated. Intially, this comprises all terms appearing in base equalities. We
then group congruent terms based on the equalities currently discovered, and
merge the congruence classes of each group in the union-find. Grouping is
accomplished via semisort---the \texttt{group\_by} operation in
\texttt{parlaylib}---and merging congruence classes is accomplished through
taking the union of the corresponding equality classes in the union-find.

We then propagate terms whose congruence classes may have changed to the next
round. These are exactly the \emph{parents} of nodes that whose congruence
classes have changed that round. Specifically, if the congruence class of term
$e$ has changed, then the congruence class of $f(e)$ may have changed as well.
This requires maintaining the parents of every congruence class; specifically a
mapping \textit{pred} between class identifiers and identifiers of parents of
that class. This mapping must be updated every round; if two congruence classes
are merged, then their parents must be merged as well. These rounds continue
until there is no longer any work to do. Pseudocode depicting this algorithm is
shown in Figure~\ref{fig:alg:bsp}.

The algorithm is decomposed into three different functions. The function
\textsc{Congruent-Level1} checks whether two terms are congruent based on the
equivalence classes in the union find\footnote{\textsc{Level1} comes from the
fact that we are only considering congruent terms with only one level of
nesting}. Two terms are congruent if they are both the same function symbol and
each of their subterms are pointwise congruent.

The \textsc{MergeCongruenceClass} function is used to merge the equivalence
classes of different terms. Specifically, given a sequence $S$ of terms, it will
union the classes of every term in $S$. We take a simple divide-and-conquer
approach; we merge the left and right halves (in parallel), and then
\textsc{Union} the representatives of the two classes.

\begin{figure}
  \centering
\begin{algorithmic}[1]
    \Function{Congruent-Level1}{X, Y} \If{$X = f(M_1,\ldots,M_n)$ and $Y =
      f(N_1,\ldots,N_n)$} \State \Return $\Call{Find}{M_i} = \Call{Find}{N_i}$
      for all $i = 1,\ldots,n$ \Else \State \Return \textbf{false} \EndIf
      \EndFunction

    \Function{MergeCongruenceClass}{\textit{Class}} \If{$\mathit{Class} =
      \{e\}$} \State \Return{$e$} \Comment{Singleton class, return
      representative} \EndIf \State $n \gets |\textit{Class}|$ \State $L \gets
      \mathit{Class}[0,\ldots,\frac{n}{2}]$ \State $R \gets
      \mathit{Class}[\frac{n}{2},\ldots,n]$ \State $\rhd$ Merge left and right
      halves, get their representatives \State $l, r \gets
      \Call{MergeCongruenceClass}{L} \parallel \Call{MergeCongruenceClass}{R}$
      \State \Return $\Call{Union}{l, r}$ \Comment{Union left and right reps,
      return root} \EndFunction

    \Function{ParallelClose}{equalities = $\{e_1 = e_1', \ldots, e_w = e_w'\}$}
        \State $\mathit{Frontier} \gets \text{every term appearing in } \mathit{equalities}$
        \While{$\mathit{Frontier}$ is not empty}
          \State $\mathit{groups} \gets \text{congruence classes quotiented by } \textsc{Congruent-Level1}$.
          \ParFor{$\mathit{group} \in \mathit{groups}$}
            \State $\rhd$ Merge every term within congruence class
            \State $\mathit{rep} \gets \Call{MergeCongruenceClass}{\mathit{group}}$
            \State $\rhd$ Construct new parent lists
            \State $\mathit{Parents}[\mathit{rep}] \gets \Call{ConcatMap}{\lambda \mathit{group}. \; \mathit{Parents}[\mathit{group}]}$
          \EndParFor
          \State $\mathit{Frontier} \gets \Call{FlatMap}{\mathit{groups}, \lambda g. \; \mathit{Parent}[\Call{Find}{g}] }$
        \EndWhile
    \EndFunction
\end{algorithmic}

\caption{Round-based congruence closure algorithm}
\label{fig:alg:bsp}
\end{figure}
